\chapter{All proofs inside one finite-dimensional space}\label{tx:ch:closure}\label{tx:sec:pc-space}
A proof can be arbitrarily long even when its degree is fixed. To study
all such proofs, we need not enumerate their possible lengths. There
are only finitely many monomials of degree at most $d$ in a fixed finite
set of variables. Every line of a degree-$d$ proof is a vector in the
space on those monomials. The proof rules describe which vectors can
be reached.

This chapter identifies the reachable space exactly. The distinction
between NS and PC becomes the distinction between one linear span and
an iterated closure. In particular, the closure must recognise a
low-degree vector created by cancellation among higher-degree vectors.

\section{Coordinates for bounded-degree polynomials}
Let $\mathcal X$ be the finite set of variable names and put
$v=|\mathcal X|$. Work over the fixed field $K$. The space $P_{\le d}$
consists of polynomials in $x_u$, $u\in\mathcal X$, of total degree at
most $d$. Its monomial basis has size
\[
                     \dim P_{\le d}=\binom{v+d}{d}.
\]
Indeed, a monomial has exponents $a_1,\ldots,a_v$ with
$\sum a_i\le d$. Add the slack coordinate
$a_{v+1}=d-\sum a_i$. This gives $v+1$ nonnegative integers summing
to $d$. Placing $v$ separators among $d+v$ positions counts those
choices by $\binom{v+d}{v}=\binom{v+d}{d}$.

For two variables and $d=2$, one possible ordering is
$1,x,y,x^2,xy,y^2$. A polynomial is its vector of coefficients in this
basis. Addition and scalar multiplication are the usual vector-space
operations; polynomial multiplication is a different operation which
may leave this space.

Fix a finite system $\Gamma$. We omit its zero inputs when forming
spanning sets, since zero contributes no new vector. For NS set
\begin{equation}\label{tx:eq:NSspan}
 I_d=\Span_K\{\mu f:f\in\Gamma,\ \mu\text{ a monomial},\quad
                              \deg\mu+\deg f\le d\}.
\end{equation}
There are only finitely many permitted pairs $(\mu,f)$. If desired,
put one column $\mu f$ in a matrix for each such pair. Deriving a target
by NS in degree $d$ will mean that its coefficient vector belongs to
the column span of this matrix.

\section{The closure for ordinary PC}
For $d\ge1$ define
\begin{equation}\label{tx:eq:closure}
 \begin{split}
 V_0&=\Span_K\{f\in\Gamma:\deg f\le d\},\\
 V_{j+1}&=V_j+\sum_{u\in\mathcal X}x_u(V_j\cap P_{\le d-1}),\\
 C_d&=V_\infty:=\bigcup_{j\ge0}V_j.
 \end{split}
\end{equation}
For $d=0$, define $C_0=V_0$ without a multiplication term.
The notation $x_uU$ means the image of a subspace $U$ under the linear
map $f\mapsto x_uf$. The sum of subspaces consists of finite sums of
vectors from them.

The intersection in this formula is crucial. It asks which vectors
already derivable in degree $d$ happen to have actual degree at most
$d-1$. It does not merely select input axioms of degree at most $d-1$,
nor basis vectors that were short when first listed. A linear combination
may cancel its highest terms and enter the intersection.

\begin{example}[The first closure step sees cancellation]
Return to $f_1=xy+x$, $f_2=xy+1$, $f_3=xz$ over $\mathbb F_2$, with
$d=2$. A combination
\[
                   af_1+bf_2+cf_3
                  =(a+b)xy+cxz+ax+b
\]
is linear precisely when $a=b$ and $c=0$. Thus
\[
                     V_0\cap P_{\le1}=\Span\{x+1\}.
\]
The vector $z(x+1)$ lies in the new multiplication space, and adding
$f_3$ produces $z\in V_1$. The degree-two NS space, in contrast, is
just the span of the three quadratic inputs and contains no $z$ term.
\end{example}

\begin{proposition}[Exact description of consequences]\label{tx:prop:closure}
For every polynomial $G$ and every $d\ge0$,
\[
 \begin{aligned}
 h_{\NS}(\Gamma\vdash G)\le d&\iff G\in I_d,\\
 h_{\PC}(\Gamma\vdash G)\le d&\iff G\in C_d.
 \end{aligned}
\]
The chain $V_j$ stabilises after finitely many strict increases, and
$I_d\subseteq C_d\subseteq P_{\le d}$.
\end{proposition}
\begin{proof}
Expand every coefficient $a_f$ in an NS identity into monomials.
Every resulting product belongs to the spanning set in
\eqref{tx:eq:NSspan}. Conversely, group a linear combination of such
products by input $f$; the resulting multiplier has degree at most
$d-\deg f$. This proves the NS equivalence, including the empty
identity for $G=0$.

For PC, first prove by induction on a derivation that every line is in
$C_d$. Usable input axioms belong to $V_0$. Two vectors in stages $V_i$
and $V_j$ both belong to $V_{\max(i,j)}$, which is closed under linear
combinations. If a nonzero line $f$ is multiplied by a variable while
staying within degree $d$, then $\deg f\le d-1$. If $f\in V_j$, its
multiple belongs to $V_{j+1}$ by definition. Multiples of zero are zero.

For the converse, induct on the closure stage. Every vector of $V_0$
is a finite linear combination of usable axioms, hence has a legal
proof. A vector of $V_{j+1}$ is a finite sum of an old vector and
terms $x_vf_v$ with $f_v\in V_j\cap P_{\le d-1}$. Concatenate proofs
of the finitely many $f_v$, append the permitted multiplications, and
add to obtain the desired vector as the last line. Every partial sum
still has degree at most $d$.

The chain is increasing inside $P_{\le d}$. A strict increase raises
its dimension, so there are at most $\dim P_{\le d}$ such increases.
If $V_{j+1}=V_j$, the recurrence gives equality at every later stage.
Thus one stage equals the whole union. Finally each NS generator
$\mu f$ has a PC derivation within the same bound by successive
variable multiplication, giving $I_d\subseteq C_d$.
\end{proof}

The proposition includes two different finiteness statements. Every
particular vector in the closure has a finite proof, by the constructive
induction. The space containing \emph{all} degree-$d$ proofs is also
reached after finitely many stages. Neither statement assumes in advance
an upper bound on the length of proofs presented to us.

\section{The linear maps behind the recurrence}
The closure formula can be expressed entirely through fixed ambient
spaces and linear maps. This is the form needed for transferring it
between universe sizes.

The immediate idea is elementary. To describe all permitted variable
multiplications, keep a separate low-degree input polynomial $f_v$
for each variable $x_v$, and send this list to $\sum_v x_vf_v$.
This is a linear map on lists. Later, in
Section~\ref{tx:sec:multiplication-registers}, we call the space of
these lists $X(n)\otimes L_d(n)$ and the map $m_n$.
The tensor notation packages the separate registers; it does not
introduce a new PC inference rule. All these operations occur at
one fixed universe size. Lifting to another size is a separate task.

The first map sends an axiom index to its coefficient vector. Restrict
to axioms of degree at most $d$; its image is $V_0$. The NS map instead
sends a permitted pair $(\mu,a)$ to $\mu f_a$; its image is $I_d$.
Both use the actual input coefficients.

Let $X$ be the vector space with one basis element for each polynomial
variable, and put $L=P_{\le d-1}$. Multiplication is represented by a
single linear map
\[
              m:X\otimes L\longrightarrow P_{\le d},
                         \qquad x_v\otimes f\longmapsto x_vf.
\]
Here the tensor product can be understood concretely: its basis is
made of pairs $(x_v,\mu)$ with $\deg\mu\le d-1$, and the map sends
such a pair to the monomial $x_v\mu$. Bilinearity is built into the
vector space on these pairs.

For any current stage $V_j$, its low-degree part is the intersection
$V_j\cap L$. It can be described as the kernel of
\[
                  V_j\oplus L\longrightarrow P_{\le d},
                              \qquad(f,g)\longmapsto f-g,
\]
identified with either equal coordinate. Multiplication then gives an
image, and adjoining old vectors gives a sum:
\[
 V_{j+1}=V_j+\im\bigl(m|_{X\otimes(V_j\cap L)}\bigr).
\]
Thus the construction uses a small vocabulary: coefficient maps,
kernels, intersections, images, sums, and a finite iteration. A
transport operation preserving this vocabulary, and agreeing with the
actual maps, transports the closure. Merely matching dimensions would
not identify which target vectors lie in it.

\section{Changing the coefficient field}
Extending scalars enlarges the allowed coefficients while keeping the
same coordinate equations. If $U\subseteq K^M$ has basis
$u_1,\ldots,u_a$ and $K\subseteq E$, its scalar extension is
$U_E=\Span_E\{u_1,\ldots,u_a\}\subseteq E^M$.
It is also written $E\otimes_K U$. The subscript $K$ records which
scalars may be moved between the two tensor factors.
For example, the kernel of $(a,b)\mapsto a+b$ is $K(1,-1)$ over
$K$ and $E(1,-1)$ over $E$. The same basis and relation remain,
although there are more choices of scalar multiples.
The next lemma applies this coordinate principle to every stage
of the proof calculation.

\begin{lemma}[Field extension preserves minimum degree]\label{tx:lem:field}
If $K\subseteq E$ are fields and $\Gamma,G$ have coefficients in $K$,
then the least NS and ordinary-PC degrees of $\Gamma\vdash G$ are the
same over $K$ and over $E$.
\end{lemma}
\begin{proof}
At a fixed degree, NS membership is solvability of a finite linear
system with coefficients in $K$. Row reduction over $K$ already
determines whether it is consistent; extending scalars does not change
that answer. Equivalently, a matrix and its augmented matrix have the
same ranks over the two fields.

For PC, extending scalars commutes with sums and images of finite
matrices, and also with their kernels: row reduction describes these
spaces over $K$ and gives bases after extension. Intersections are
kernels of the difference map just displayed. Induction on
\eqref{tx:eq:closure} therefore identifies each extended stage with the
stage computed over $E$. Once both chains have stabilised, membership
of the original vector $G$ is unchanged. This proves equivalence at
every integer degree, including degree zero, and hence equality of
the minimum degrees, including infinity.
\end{proof}

This lemma will let us use convenient auxiliary coefficients in the
finite-dimensional construction. It does not change the target to an
arbitrary polynomial over the larger field. The target whose membership
is tested remains the original $K$-coefficient vector.

\section{When degree gives a proof-size bound}
This section is a separate application of the consequence-space view.
It explains the size interpretation under full coefficient lists; the
construction of the lifting operator and the degree lower bounds do
not depend on this encoding. A reader following the degree proof can
continue to Chapter~\ref{tx:chap:contract} and return here for the size
consequences.

Degree alone is not a count of the nonzero monomials in a proof. We use
one particular, deliberately explicit representation in which a degree
bound has a direct size interpretation.

\begin{definition}[Full coefficient-list size]\label{tx:def:dense-size}
In an ambient ring with $v$ variables, store a nonzero polynomial of
degree $e$ by the full list of coefficients of all monomials of degree
at most $e$, including zeros. The number of entries is
$B_v(e)=\binom{v+e}{e}$. Store zero by one entry. PC size is the sum of
these lengths over all lines, including the final target. NS size is
the sum over its nonzero multipliers; the inputs and target are supplied
separately. The empty NS certificate for zero has size zero.
\end{definition}

Write $S_{\NS}$ and $S_{\PC}$ for the least coefficient-list sizes of
proofs of the fixed target $G$ from $\Gamma$. These are minimum sizes
over all proofs in the specified system; they are infinity when no
proof exists.

\begin{lemma}[Degree versus coefficient lists]\label{tx:lem:dense-size}
With $v$ variables and input degree at most $D$, a degree-$d$ NS proof
has a representative certificate of size at most $B_v(d)^2$. A
 degree-$d$ PC proof has a representative proof of size $O(B_v(d)^3)$.
Conversely, if the least degrees are finite, minimum sizes satisfy
\[
 S_{\PC}\ge B_v(h_{\PC}),\qquad
 S_{\NS}\ge B_v(\max\{0,h_{\NS}-D\})\quad(G\ne0).
\]
Furthermore $B_v(d)\ge2^{\min(v,d)}$.
\end{lemma}
\begin{proof}
Put $M=B_v(d)$. A vector in $I_d$ is a linear combination of at most
$M$ generator columns. Grouping by input gives at most $M$ nonzero
multipliers, each representable by at most $M$ entries.

For PC retain a basis of usable axioms. If its span $W$ is not closed
under the permitted multiplications, choose
$f\in W\cap P_{\le d-1}$ with $x_vf\notin W$. Derive $f$ as a
combination of the at most $M$ retained basis vectors, append $x_vf$,
and enlarge the basis. Each enlargement raises dimension, so there are
at most $M$ enlargements. At the end, the closed space contains every
degree-$d$ consequence. Append a combination giving the desired target.
This uses $O(M^2)$ lines, each with at most $M$ entries.

Any PC proof has a line of degree at least the least possible proof
degree. An NS certificate for a nonzero target has a nonzero product
of degree at least its least degree; its multiplier has degree at least
$h_{\NS}-D$. These give the lower bounds. Finally the table contains
all squarefree monomials in any $\min(v,d)$ variables, of which there
are $2^{\min(v,d)}$.
\end{proof}

All variables of the ambient ring count, including those absent from
nonzero terms. Over a fixed finite field, coefficient entries have
bounded bit length. The results do not automatically become lower
bounds for sparse representations, where zero coefficients are omitted.
This distinction will remain explicit in the size classification.

\section{Exercises}
\begin{exercise}\label{tx:ex:closure-zero}
Describe $I_0$ and $C_0$ for an arbitrary system. Prove directly that a
nonzero constant target has a degree-zero proof exactly when an input
is a nonzero constant.
\end{exercise}
\noindent{\small\hyperref[tx:ans:closure-zero]{Model answer (p.~\pageref*{tx:ans:closure-zero}).}}\par\medskip

\begin{exercise}\label{tx:ex:closure-witness}
With input $\Gamma=\{x^2\}$, show that every vector in every closure
$C_d$ has zero coefficient of $x$. Deduce that $x$ is underivable,
although it vanishes at every field-valued solution of $\Gamma$.
\end{exercise}
\noindent{\small\hyperref[tx:ans:closure-witness]{Model answer (p.~\pageref*{tx:ans:closure-witness}).}}\par\medskip

\begin{exercise}\label{tx:ex:dense-count}
List the monomials counted by $B_2(2)$, and compare the full coefficient
list of $x^2$ with its sparse representation. For $v\ge d$, identify
explicitly the $2^d$ entries used in the final lower bound of
Lemma~\ref{tx:lem:dense-size}.
\end{exercise}
\noindent{\small\hyperref[tx:ans:dense-count]{Model answer (p.~\pageref*{tx:ans:dense-count}).}}\par\medskip

